\section{The finite-jet obstruction}\label{sec:finite-jet}
We now convert the exclusion of a formal critical arc into a finite analytic estimate. This step uses an established formal real Nullstellensatz and ordinary analytic approximation, each in finitely many variables.

Let $(b,x)$ be the stationary coordinates of \Cref{lin:family}, and write the formal action as
\[
 \mathcal F(b,x)\in\R\{b\}[[x]].
\]
It is a formal series in the opening variables $x$, with coefficients convergent in $b$. For $N\geq0$, let $f_N$ denote its Taylor polynomial through total $x$-degree $N$. Gradients below include both the $b$ and $x$ coordinates.

\begin{lemma}[Finite-power estimate]\label{jet:finite-power}
There is an integer $m\geq1$ such that, for every sufficiently large fixed $N$, there are a neighborhood of $(0,0)$ and a constant $C_N$ for which
\begin{equation}\label{jet:inequality}
 |x|^m\leq C_N\bigl(|\nabla f_N(b,x)|+|x|^{N-1}\bigr)
\end{equation}
for real $(b,x)$ in that neighborhood. The exponent $m$ is independent of $N$. If there are no cylindrical ends, there are no $x$ variables and the assertion is vacuous.
\end{lemma}

\begin{proof}
Regard $\mathcal F$ as an element of the formal ring $\R[[b,x]]$, and let $I$ be the ideal generated by its finitely many first partial derivatives. By \Cref{phase:no-arc}, every formal arc through the origin on which $I$ vanishes has $x=0$. The formal real arc criterion of Aschenbrenner--Srhir \cite[Corollary~4.8]{AS2024} therefore puts each $x_i$ in the real radical of $I$. In particular, there are integers $m_i\geq1$ and formal series $h_{i\nu},a_{i\ell}$ such that
\begin{equation}\label{jet:certificate}
 x_i^{2m_i}+\sum_\nu h_{i\nu}(b,x)^2
   =\sum_\ell a_{i\ell}(b,x)\,\partial_\ell\mathcal F(b,x).
\end{equation}
All sums in \eqref{jet:certificate} are finite. These certificates, and hence the integers $m_i$, are chosen once.

Fix $N>2\max_i m_i+1$ and reduce \eqref{jet:certificate} modulo the ideal $(x)^{N-1}$. Only finitely many coefficients in $x$ of the unknown series occur. They are formal functions of the same variables $b$. The coefficients of $\mathcal F$ which occur are convergent functions of $b$, so coefficient comparison gives a finite system of convergent analytic equations in $b$ and those finitely many unknown functions. The formal coefficients of \eqref{jet:certificate} solve that system. After translating their constant terms to the origin, Artin's analytic approximation theorem \cite[Theorem~1.2]{Artin1968} gives convergent coefficient functions solving it exactly. No nested approximation requirement is imposed, and the infinite $x$-series is not being approximated all at once.

Let $H_{i\nu}$ and $A_{i\ell}$ be the resulting polynomials in $x$ with convergent coefficients in $b$. Replacing the gradient of $\mathcal F$ by that of $f_N$ changes nothing modulo $(x)^{N-1}$. Consequently, for convergent real functions near the origin,
\begin{equation}\label{jet:analytic-certificate}
 x_i^{2m_i}+\sum_\nu H_{i\nu}(b,x)^2
 =\sum_\ell A_{i\ell}(b,x)\,\partial_\ell f_N(b,x)
       +O_N(|x|^{N-1}).
\end{equation}
The remainder estimate is uniform for $b$ in a sufficiently small fixed neighborhood: every monomial in its finite polynomial expression has $x$-degree at least $N-1$, and its coefficient is bounded there. Evaluation at real points makes the sum of squares nonnegative, and the coefficients $A_{i\ell}$ are bounded. Thus
\[
 |x_i|^{2m_i}\leq C_N\bigl(|\nabla f_N|+|x|^{N-1}\bigr).
\]
Set $m=2\max_i m_i$. For $|x|\leq1$, each $|x_i|^m\leq|x_i|^{2m_i}$; summing and using equivalence of finite-dimensional norms proves \eqref{jet:inequality}.
\end{proof}

The order of choices in \Cref{jet:finite-power} matters. The radical certificates fix $m$ first. One then chooses a single sufficiently large $N$ and applies the ordinary analytic gradient inequality to that one function $f_N$. The spatial localization loss may subsequently tend to zero without changing the gradient exponent. This is the order used in \Cref{sec:flow}.
